\section{Implications of the Findings}
\label{sec:discussion}
\paragraph{A class of error that current evaluations do not measure.}
Error taxonomies for language models ask whether an answer contradicts the world or contradicts the input \citep{space}. Silent misreadings do neither. The input is consistent with the answer but incomplete relative to the document it came from, and the answer is wrong only with respect to that document. Such errors become visible only when an answer is compared with the full source rather than with the supplied input, and the small effect on public benchmarks indicates that this comparison is rarely made rather than that the errors are rare. Evaluations of systems that read local text, including retrieval-augmented generation, need questions that reach the boundaries of local definitions.

\paragraph{Knowing that a definition is needed is a distinct capability.}
Models largely know what they know \citep{know-what-know}, and calibration and selective prediction exploit this to abstain when uncertain \citep{calibration,selective}. A redefined common word defeats this mechanism, because the model is not uncertain about the word. It is confident in the wrong sense. The strongest readers follow a stated definition on nearly every item and misread largely the same occurrences as other readers, so the remaining failure is not a lack of capacity to use local meaning. It is the absence of any internal signal that local meaning is needed. Recognizing that a familiar word may carry a local meaning is thus a capability that neither scale nor calibration has produced, and a natural target for evaluation and training.

\paragraph{Designing systems that read local text.}
Three design consequences follow. Local definitions should be taken from the source rather than generated by the model. Safeguards based on confidence or sample agreement, and faithfulness-oriented decoding, should not be expected to catch these errors. In documents whose definitions refer to other material, retrieval must follow those references. The same considerations apply wherever text is written for a group whose vocabulary differs from general usage, as in software engineering terminology \citep{se-terminology}, community slang \citep{informal-language}, and organizational documents. Whether item-level difficulty rankings such as $R$ transfer to requests written by users is a direction for future work.
